% Pursuing Stacks (cotes 134-2 to 134-8): archive cote 134-6, site page 90 (PDF page 91), typescript page 408.
% typescript. Draft compared: ps5.tex:6066-6142.
% AI-transcribed by Claude Opus 5.5 (claude-opus-5-5, adaptive thinking, effort high), Message Batches API, batch ps3 (msgbatch_01XdFmuVVBAcgc7EFxdtCqsj), 2026-11-12; reviewed against both supplied scan halves; see CHECK_DECISIONS.tsv.
\scanpage{134-6/90}
\mspage{408}

lent to the four others,  namely \xout{\uline{Homaddinf(A}}
\[ \underline{\text{Homaddinf}}(\text{Addinf}(A),\hand{\mathcal{M}}) \quad , \]
the category of infinitely additive functors from the infinitely additive
envelope of A into \hand{$\mathcal{M}$} (cf.  section 99 p.366 for description of the cate-%
gory Addinf(A)).  Rather than writing down the larger diagram here,  I'll
write down the dual enlarged one\hand{,} for the \xout{topos} dual topos B\^{} and \inserted{for} various
expressions of the category of \hand{\uline{\textcolor{black}{sheaves}}} on this topos,  with values in a
category \hand{$\mathcal{N}$} stable this time under small inverse limits:
\diagram{
(62)\quad
% Hand-drawn arrows are marked blue; typed objects remain plain.
\begin{tikzcd}[column sep=small]
 & \underline{\text{Hom}}^{\hand{!}}(B^{\wedge o},\hand{\mathcal{N}}) \arrow[dr,color=blue,"\hand{\approx}"] & \\
\underline{\text{Hom}}^{\hand{!}}(B^{\wedge}_{ab}{}^{o},\hand{\mathcal{N}}) \arrow[ur,color=blue,"\hand{\approx}"] \arrow[d,color=blue,"\hand{\approx}"'] & & \underline{\text{Hom}}(B^{o},\hand{\mathcal{N}}) \\
\underline{\text{Hommultinf}}(\text{Addin}(B)^{o},\hand{\mathcal{N}}) \arrow[r,color=blue,"\hand{\approx}"] & \underline{\text{Homadd}}(\text{Add}(B)^{o},\hand{\mathcal{N}}) \arrow[ur,color=blue,"\hand{\approx}"] &
\end{tikzcd}
\quad ,
}

\over[typed]{w}{W}here \uline{Hommultinf}  denotes the category of``infinitely multiplicative''
functors from one additive category stable under infinite products to
an\over{\unsure{oh}}{ot}her.  Recalling for the extreme right term of (62) that B$^{o}$=A,  we see
that this term is \hand{\uline{\textcolor{black}{identic}}} to the corresponding term in the diagram (even
the enlarged one) (58)  -  hence,  if \hand{$\mathcal{M}$} is stable under both types of
limits,  the t\over[typed]{\illegible{1}}{e}n categories occuring \inserted[typed]{altogether} in the two diagrams are mutually
equivalent (as a matter of fact,  there are nine only which are mutually
different).  It may be noted that there is still another pair of corres-%
ponding terms in the two diagrams for which the equivalence between them
may be viewed as tautological,  namely
\[ \underline{\text{Homadd}}(\text{Add}(A),\hand{\mathcal{M}}) \hand{\approxeq} \underline{\text{Homadd}}(\text{Add}(B)^{o},\hand{\mathcal{M}}) \quad , \]
due to the tautological equivalence of categories
\[ \text{Add}(B)^{o} \hand{\approxeq} \text{Add}(B^{o}) \;=\; \text{Add}(A) \qquad . \]

As emphasize\over[typed]{\unsure{e}}{d} in yesterday\hand{'}s notes,  the \over[typed]{\illegible{1}}{c}\over{\illegible{1}}{a}nonical pairing (deduced
from (59)  by taking \hand{$\mathcal{M}$}= (Ab))
\[ (63)\qquad A^{\wedge}_{ab} \text{ x } B^{\wedge}_{ab} \hand{\longrightarrow} (\text{Ab}) \quad , \qquad (F,G) \hand{\longmapsto} F *_{\underline{Z}} G \hand{,} \]
deserves special attention,  giving rise to an equivalence between each
of the mutually dual \inserted[typed]{abelian} topoi A$^{\wedge}_{ab}$ , B$^{\wedge}_{ab}$  with the category of abelian coshea-%
ves on the other
\[ (64)\qquad B^{\wedge}_{ab} \hand{\xrightarrow{\approx}} \underline{\text{Hom}}_{\hand{!}}(A^{\wedge}_{ab},(\text{Ab})) \quad , \quad A^{\wedge}_{ab} \hand{\xrightarrow{\approx}} \underline{\text{Hom}}_{\hand{!}}(B^{\wedge}_{ab},(\text{Ab})) \]
(cf (49) and (49\hand{'})  page 399).  It should be kept in mind that besides
this duality pairing between two abelian topoi,  there is important extra

\margin{90}
